\documentclass[11pt]{amsart}
\usepackage[T1]{fontenc}
\usepackage{lmodern}
\usepackage{amssymb}
\usepackage[hidelinks]{hyperref}
\newtheorem{theorem}{Theorem}[section]
\newtheorem{proposition}[theorem]{Proposition}
\newtheorem{lemma}[theorem]{Lemma}
\newtheorem{fact}{Fact}
\renewcommand{\thefact}{F\arabic{fact}}
\numberwithin{equation}{section}
\DeclareMathOperator{\rad}{rad}
\newcommand{\eps}{\varepsilon}
\newcommand{\Z}{\mathbb Z}
\newcommand{\Q}{\mathbb Q}
\newcommand{\R}{\mathbb R}
\newcommand{\parameters}{%
\begin{table}[htbp]\centering\small
\caption{Parameters used in the proof}\label{tab:parameters}
\begin{tabular}{@{}l@{\quad}p{.73\linewidth}@{}}
\hline
Symbol & Choice or constraint\\
\hline
$k$ & $20$ for the polynomial bound.\\
$\eta,\lambda,\nu$ & $0<\eta<1/9$ rational, $\lambda=1-3\eta$,
 $\nu=(\lambda-2/3)/100$.\\
$\sigma,\kappa$ & $0<\sigma\le\nu$, $0<\kappa<(1-9\eta)/16$.
 The exterior estimates use $\sigma=\nu$, $\kappa=(1-9\eta)/32$.\\
$Y,r,M$ & \raggedright $Y=H\exp(-\sigma x)$, $r=2x/\log x$ for $x>1$,
 $M=\max(1,p/H)\exp(\nu x)$.\tabularnewline
$b_0,a_0$ & $b_0=\nu/2$, $a_0=3\eta+(\lambda-2/3)/2$.\\
$\theta,Z$ & $0<\theta<1$, $3a_0(1+\theta)^2<1-2\kappa$.
 Choose $0<Z<1$ satisfying \eqref{eq:slack}.\\
$k_0,W$ & $k_0=4/(3Z)$, $W=\lceil4/(3Zb_0)\rceil+1$.
 The excluded denominators are $1\le Q\le W$.\\
$\alpha,\rho,\delta$ & $\alpha=3(1+\theta)/2$, $\rho=\alpha ht/\ell$,
 $\delta=Z\rho$.\\
$A,N_0,K,N$ & $A=2/(ht)$, $N_0=A\rho^2$,
 $K=\lfloor((1-Z)^2N_0+1)/2\rfloor$, $N=2K$.\\
$\rho_1,B$ & $\rho_1=\rho+\delta$, $B=2^{\rho_1/t+K}$.\\
$X$ & Lower bound for $x$. At assembly take the maximum of the thresholds
 for the same $\eta,\sigma,\kappa$.\\
\hline
\end{tabular}
\end{table}}

\title[An upper bound for log c]{An upper bound for log c in the abc conjecture with exponent 2/9}
\author{Ruiyi Zhang}
\address{Independent Researcher, Columbia, Maryland, USA}
\email{ryzfsu@gmail.com}
\date{November 20, 2026}
\subjclass[2020]{Primary 11D75; Secondary 11J86}
\keywords{abc conjecture, radical, logarithmic forms, interpolation determinant}
\hypersetup{
  pdfauthor={Ruiyi Zhang},
  pdftitle={An upper bound for log c in the abc conjecture with exponent 2/9},
  pdfsubject={MSC 2020: Primary 11D75, Secondary 11J86},
  pdfkeywords={abc conjecture, radical, logarithmic forms, interpolation determinant}
}
\begin{document}
\begin{abstract}
Put $R=\rad(abc)$ for positive coprime integers $a+b=c$.
We obtain $\log c\ll_\eps R^{2/9+\eps}$ for every $\eps>0$.
A large prime power on the side of least radical determines one rational
number close to $1$ at its prime. Two auxiliary primes give a second,
multiplicatively independent number. On their plane a determinant has
greater divisibility than its row heights allow.
\end{abstract}
\maketitle
\section{Introduction}
Put $R=\rad(abc)=\prod_{p\mid abc}p$. The $abc$ conjecture predicts
$c\ll_\eps R^{1+\eps}$ for positive coprime integers $a+b=c$.
The conjecture remains open. Stewart and Yu obtained
$\log c\ll R^{1/3}(\log R)^3$ \cite[Theorem~1]{SY2001}.
The bound in \cite[Theorem~1, (1.1)]{Zhang2026Third} is
$\log c\ll R^{1/3}\exp(-\log R/(3000\log\log R))$ for large $R$.

\begin{theorem}\label{thm:main}
For every real $\eps>0$ there exists $C_\eps>0$ such that
\[
                 \log c\le C_\eps R^{2/9+\eps}
\]
for all coprime positive integers $a,b$ and $c=a+b$.
\end{theorem}

Write $H=\log c$. Equation (1.2) of \cite{Zhang2026Third} uses
$W=\log R$ and $X=R^{1/3}$; here we put $x=W/3$.
Order the summands as $a\le b$.
For each $\eta$ as in Section~\ref{sec:cap} we exclude, for large $x$, the range
\begin{equation}\label{eq:high}
                     H\ge R^{1/3-\eta}=\exp((1-3\eta)x).
\end{equation}
The polynomial bound of Section~\ref{sec:cap} then suffices on bounded $x$.

Let $s$ be a side of least radical. Sections~\ref{sec:real} and
\ref{sec:small} show that $\log s$ is close to $H$ and that small
primes contribute little to it. Since $\rad(s)\le\exp(x)$,
one prime $p$ carries almost all of $H$ in $v_p(s)\log p$
(Section~\ref{sec:sel}).

The signed ratio of the other sides is $1$ modulo $p^m$, where $m=v_p(s)$.
Section~\ref{sec:pair} finds an independent rational number congruent to
$1$ modulo $p$, using two new primes. Put $\ell=\log p$.
On the plane of their exponents, the height is bounded by
$g(i,j)=h|i+\xi j|+t|j|$, with $H-\log2\le h\le H$.

Put $\rho=\alpha ht/\ell$. The diamond $g\le\rho$ has area
$N_0=2\rho^2/(ht)$ and its mean
height bound is $2\rho/3$. Section~\ref{sec:determinant} constructs a nonzero
minor of order $N=2K$ from this diamond and a short translate.
When $N\le3m/2$, the divisibility estimate in that section gives
\[
  v_p(\Delta)\ge
  \min_{0\le y\le K,\; y\in\mathbb R}
    \big((K-y)4K/3+y^2\big)=8K^2/9=2N^2/9.
\]
The minimum occurs at $y=2K/3$. At the level of areas,
$2N_0\rho/3<2N_0^2\ell/9$ requires $\alpha>3/2$.
The depth budget at $m=h/\ell$ requires $\alpha^2t/\ell\le3/4$.
These inequalities allow slack precisely when $t/\ell<1/3$.
Section~\ref{sec:determinant} obtains $t/\ell<a_0$ with
$3\eta<a_0<1/3$. Such an interval exists precisely when $\eta<1/9$.
Section~\ref{sec:end}
lets $\eta$ approach $1/9$, giving $1/3-1/9=2/9$.

The use of AI tools is described at the end of the paper.

\section{A polynomial bound}\label{sec:cap}
\begin{proposition}\label{lem:cap}
There are real numbers $C_0>0$ and $k\ge0$ such that $H\le C_0R^k$
for every triple.
\end{proposition}
\begin{fact}[{\cite[Fact F1, (2.1)]{Zhang2026Third}}]\label{fact:sy}
For $c>2$ one has $H\le C_1 R^{1/3}(\log R)^3$ with absolute $C_1$.
\end{fact}
\begin{proof}[Proof of Proposition~\ref{lem:cap}]
Apply Fact~\ref{fact:sy} with the radical of the three pairwise coprime sides.
Since $\log R\le R$, $k=20$ and $C_0=\max(1,C_1)$ suffice, including $(1,1,2)$.
\end{proof}

Fix a rational $\eta$ in $(0,1/9)$ and use Table~\ref{tab:parameters}.
Choose the first of $a,b,c$ having least radical and denote it by $s$.
Define $z=c/b,c/a,-a/b$, respectively, when $s=a,b,c$.
These are the ratios in \cite[(3.4)]{Zhang2026Third}, inverted in the first two cases.
For $q\mid s$ they are $q$-units, so inversion preserves $v_q(1-z)=v_q(s)$.
Use the normalizations of \cite[\S2]{Zhang2026Third}, with $v_q(q)=1$.
For $n/d$ in lowest terms with $d>0$ put
$h(n/d)=\log\max(|n|,d)$ and write $h=h(z)$.
Set
\[
  S(s,Y)=\sum_{\substack{q\mid s,\ q\le Y\\q\ {\rm prime}}}v_q(s)\log q .
\]
For the prime $p$ of Section~\ref{sec:sel} put $\ell=\log p$ and $m=v_p(s)$.
The length $t$ in Table~\ref{tab:parameters} is defined in
\eqref{eq:plane-height}. All logarithms are natural.
\parameters

\section{The real estimate}\label{sec:real}
\begin{proposition}\label{lem:real}
For every $\kappa>0$ there is $X\ge0$
such that $\log(c/a)\le\kappa H$ for every ordered triple with $x\ge X$
satisfying \eqref{eq:high}.
\end{proposition}

\begin{fact}\label{fact:local}
For every $\vartheta>0$ there is $X_{\vartheta}\ge0$ with the following
property. If $x\ge X_{\vartheta}$, then for every signed quotient
$u\ne\pm1$ of two sides and every prime $q\le R$,
\[
 -\log|1-u|\le \exp(\vartheta x),\qquad
 v_q(1-u)\log q\le q\exp(\vartheta x).
\]
\end{fact}
\begin{proof}
Apply \cite[Fact F3, Lemma~2.1, (2.3)--(2.4)]{Zhang2026Third}:
$u$ has height at most $H$ and lies in the group generated by $-1$
and the nonempty prime support of $R$. Use (2.6) with $y=1,q\le R$
and (2.8) with $J=R$, $\gamma=\vartheta/6$.
Since $1/\log q<2$, the common factor is at most
$2C_\gamma C_H(1+3x)\exp(\vartheta x/2)\le\exp(\vartheta x)$
for large $x$.
\end{proof}
\begin{proof}[Proof of Proposition~\ref{lem:real}]
Take $u=b/c$ and $\vartheta=\lambda/2$ in Fact~\ref{fact:local}.
By \eqref{eq:high}, $\log(c/a)/H\le\exp(-\lambda x/2)\le\kappa$ for large $x$.
\end{proof}

\section{The small prime powers}\label{sec:small}
\begin{proposition}\label{lem:small}
For each $\eta$ there are admissible $\sigma,\kappa$ and a threshold
$X\ge0$ such that every ordered triple with $x\ge X$ satisfying
\eqref{eq:high} has
\[
 S(s,H\exp(-\sigma x))\le\kappa H .
\]
\end{proposition}
\begin{proof}
Choose the values in Table~\ref{tab:parameters}. For $x\ge1$ the ratio
$z$ is different from $\pm1$. Since $v_q(1-z)=v_q(s)$ for $q\mid s$,
Fact~\ref{fact:local} with $\vartheta=\sigma/2$ gives $v_q(s)\log q\le q\exp(\sigma x/2)$
simultaneously beyond one threshold. There are at most $3x/\log 2$
prime divisors of $R$. Thus
\[
    S(s,Y)/H\le (3x/\log2)\exp(-\sigma x/2)\le\kappa
\]
for all sufficiently large $x$.
\end{proof}

\section{Selection of a prime}\label{sec:sel}
For any admissible $\eta,\sigma,\kappa$, write $\mathcal B_X$ for the
ordered triples satisfying \eqref{eq:high} with $x\ge X$ and
\[
            \log(c/a)\le\kappa H,\qquad S(s,Y)\le\kappa H .
\]
\begin{proposition}\label{lem:sel}
For $X\ge0$, every triple in $\mathcal B_X$ has a unique prime $p\mid s$
with $p>Y$. This prime satisfies $(1-2\kappa)H\le m\ell$.
\end{proposition}
\begin{proof}
The primes above $Y$ contribute at least
$\log s-S(s,Y)\ge(1-2\kappa)H>0$, since $s\ge a$.
Two such primes would give $\rad(s)>Y^2\ge\exp(2(\lambda-\sigma)x)>\exp(x)$.
But $s$ has least radical and $\rad(a)\rad(b)\rad(c)=R$.
Thus $\rad(s)\le R^{1/3}=\exp(x)$.
\end{proof}

\section{Two auxiliary primes}\label{sec:pair}
For an odd prime $p$ and a finite set $T$ of other primes set
\[
 \beta(v)=\prod_{q\in T}q^{v_q},\qquad
 \Lambda(T)=\{v\in\Z^T:\beta(v)\equiv\pm1\bmod p\}.
\]
For $v\in\Lambda(T)$ let $\chi(v)$ be the unique choice of
$\pm\beta(v)$ congruent to $1\bmod p$. It is an injective homomorphism
by unique factorization. Use the norm $\|v\|_w=\sum_{q\in T}|v_q|\log q$,
also on coordinate subspaces. For $v_j\ne0$ write
$\pi w=(w_q-w_jv_q/v_j)_{q\ne j}$.

\begin{proposition}\label{lem:separate}
For each admissible $\eta,\sigma,\kappa$ there exists $X>0$ such that
for every triple in $\mathcal B_X$, the prime of Proposition~\ref{lem:sel} satisfies
\[
 p\ge5,\qquad (\lambda-6\nu)x\le\ell\le x,\qquad H>x^2,\qquad \log x\ge10^9.
\]
Let $T$ be the prime support of $z$. Some $v\in\Lambda(T)$ and
$j\in T$ satisfy
\[
 \chi(v)=z,\quad |v_j|=\max_{q\in T}|v_q|>0,\quad H\le3x|v_j|,
 \quad \Lambda(T)\cap\R v=\Z v.
\]
For $w\in\Lambda(T)$ either $\pi w=0$ or $\|\pi w\|_w\ge r$.
\end{proposition}
\begin{proof}
Since $p>Y\ge\exp((\lambda-\sigma)x)$ and $p\le\rad(s)\le\exp(x)$,
we have $(\lambda-\sigma)x<\ell\le x$ and $\ell\ge x/2$.
Together with \eqref{eq:high} these give the stated bounds for large $x$.
For $v_q=v_q(z)$ we have $\chi(v)=z$ since $v_p(1-z)=m$.
The absolute numerator and denominator of $z$ are the other two sides.
Their logarithms sum to $V:=\|v\|_w$, so $2(1-\kappa)H\le V\le2H$.
Since $\log a\ge(1-\kappa)H>0$, both sides exceed one and $T$ is nonempty.
As $\sum_{q\in T}\log q\le3x$, a largest coordinate satisfies $H\le3x|v_j|$.

If $v=Dv'$ in $\Lambda(T)$ with $D\ge2$, then $u=\chi(v')$ has
$u^D=z$ and $\log2\le h(u)=h/D$. Lifting the exponent gives
$m=v_p(1-u)+v_p(D)$. The numerator and denominator of $u$ differ by at
most $2\exp(h/D)$, so $m\ell\le h/D+\log2+\log D\le H/2+O(\log H)$,
contradicting $m\ell\ge(1-2\kappa)H$ for large $x$. Thus $v$ is primitive.

Suppose $0<\|\pi w\|_w<r$. Let $L_1,L_2$ be the successive minima
of $\Z v+\Z w$. In $(v,w)$ coordinates the norm ball of radius one has area
\begin{equation}\label{eq:minima}
 \frac2{V\|\pi w\|_w}\le\mathcal A\le\frac4{L_1L_2}.
\end{equation}
It contains the diamond with vertices $\pm v/V$ and
$\pm(w-(w_j/v_j)v)/\|\pi w\|_w$. The upper bound is Minkowski's
second theorem at covolume one \cite[p.~2, (2)]{Malikiosis2010}.
A shortest vector $w_1$ extends to a basis $w_1,w_*$. Write a shortest
independent vector as $d_1w_1+d_2w_*$. If $|d_2|=1$, it completes the
basis. Otherwise take $n$ nearest to $d_1/d_2$ to obtain a basis vector
$w_*+nw_1$ of norm at most $L_2/|d_2|+L_1/2\le L_2$.

Write $v=c_1w_1+c_2w_2$ in this basis of lengths $L_1,L_2$.
The diamond on $v/V,w_{3-i}/L_{3-i}$ has area
$2|c_i|/(VL_{3-i})\le\mathcal A$, so $|c_i|L_i\le2V$.
Each $L_i\ge\ell-\log2\ge\ell/2$, since the unequal numerator and
denominator of $\chi(w_i)$ are congruent modulo $p$.
If $c_i=0$, primitivity gives $L_{3-i}=V$, so \eqref{eq:minima}
would give $L_i<2r<\ell/2$.

Put $b_i=|c_i|>0$, $B_i=2L_i$ and
$b'=b_1/B_2+b_2/B_1\le2V/(L_1L_2)\le8V/\ell^2$.
Invert the $\chi(w_i)$ as needed so that $z=\alpha_1^{b_1}/\alpha_2^{b_2}$.
These are independent rational principal units, with local degree and
residue order one, $B_i\ge\max(h(\alpha_i),\ell)$ and
$v_p(\alpha_1^{b_1}-\alpha_2^{b_2})=m$. Bugeaud and Laurent
\cite[Corollary~2, p.~315]{BL96} give
\[
 m\le\frac{24pB_1B_2}{(p-1)\ell^4}
       \max\{\log(b'\ell)+0.4,10\ell,10\}^2 .
\]
The maximum is $10\ell$ since $b'\ell\le16H/\ell$ and
$\log H<\ell+\sigma x$. Using $p\ge5$, \eqref{eq:minima} and $V\le2H$ gives
\[
 Hx/4\le(1-2\kappa)H\ell\le m\ell^2\le3000B_1B_2
 <48{,}000\,Hr=96{,}000\,Hx/\log x ,
\]
contrary to $\log x\ge10^9$.
\end{proof}

For Lemma~\ref{lem:fresh}, let $p$ be an odd prime divisor of the integer $R>0$.
Let $T$ be a nonempty subset of the remaining prime divisors of $R$.
As above $3x=\log R$. Suppose $H>0$, $\log H\ge\lambda x$,
$(\lambda-6\nu)x\le\ell\le x$. Fix $v\in\Lambda(T)$, $v_j\ne0$ and $h_0>0$.
Let $V_0$ be the volume of $\sum_{q\ne j}|y_q|\log q<r/2$ in the
coordinates $(y_q)_{q\in T\setminus\{j\}}$, with $V_0=1$ in dimension zero.

\begin{lemma}\label{lem:fresh}
For each $k_0\ge0$ and integer $W\ge0$ there is a threshold
$X=X(\eta,k_0,W)\ge0$. Suppose $x\ge X$ and
\[
 H\le3x|v_j|,\quad |v_j|\log2\le h_0,\quad
 V_0\ge\exp(-7x\log(3\log x)/\log x).
\]
Assume the line and projection conclusions of Proposition~\ref{lem:separate}.
There exist primes $\sqrt M<q_1<q_2\le M$ with $q_1q_2$ coprime to
$2R$, and $u\in\Lambda(T\cup\{q_1,q_2\})$ having new coordinates $1,-1$.
Its old coordinates $u_0$ and $\xi=u_j/v_j$ satisfy
\[
 \|\pi u_0\|_w<r,\qquad
 h_0|Q\xi-P|\ge k_0\ell\quad(P,Q\in\Z,\ 1\le Q\le W).
\]
The intersection of the enlarged lattice with the real plane spanned
by $v,u$ is $\Z v+\Z u$. Its character is injective and extends the
given character on $v$.
\end{lemma}
\begin{proof}
For $p>H$ we have $\log M\le\ell-(\lambda-\nu)x<\ell/2$,
and for $p\le H$ we have $\log M=\nu x<\ell/2$.
In both cases $\log M\ge\nu x$ and $MH/p\ge\exp(\nu x)$, with $\nu$ fixed by $\eta$.
Rosser and Schoenfeld \cite[Corollary~1, (3.5), p.~69]{RS62} give
more than $M/\log M$ primes up to $M\ge17$. Removing those up to $\sqrt M$
and the $O(x)$ prime divisors of $2R$ leaves $n\ge M/(4\log M)$ eligible primes.

Set $G=\mathbb F_p^\times/\{\pm1\}$ and let $G_T$ be generated by the classes
of $T$. The index of $\Lambda(T)$ is $|G_T|$. Extend $v$ to a basis of
$\Lambda(T)$ and project the other columns along $v$.
Expansion along coordinate $j$ gives $\pi\Lambda(T)$ covolume $|G_T|/|v_j|$.
Attach one quotient torus to each coset of $G_T$ in $G$, of total volume
$\mu=(p-1)/(2|v_j|)$. Fix a representative $g_C$ of each coset $C$.
For each eligible $q$ in $C$ take $a_q\in\Z^T$ with $[\beta(a_q)]=[q]/g_C$.
The open ball of radius $r/2$ centered at the class of $\pi a_q$ embeds by
projection separation. For their multiplicity $f$ and set $E$ of overlapping
pairs, $\int f=nV_0$, so Cauchy--Schwarz gives
\[
 2|E|V_0\ge\int f(f-1)\ge(nV_0)^2/\mu-nV_0 .
\]
Since $|v_j|\ge H/(3x)$ and $V_0\ge\exp(-o(x))$, both $n$ and
$nV_0/\mu$ exceed $\exp(\nu x-o(x))$. Thus $|E|\ge\exp(2\nu x-o(x))$.

For each overlapping pair $q_1<q_2$ choose $w\in\Lambda(T)$ so that
$u_0=a_{q_2}-a_{q_1}+w$ has $\|\pi u_0\|_w<r$.
Its residue class cancels $[q_1]/[q_2]$, so appending $1,-1$ gives $u$
in the enlarged lattice. The lift need not be unique modulo $\Z v$;
fix one per pair. Adding $iv$ replaces $\xi$ by $\xi+i$ and $P$ by $P+Qi$.

For each pair failing the slope test choose a witness $P,Q$ and put
$D=Qu_0-Pv$. Then $|D_j|<k_0\ell/\log2\le k_0x/\log2$.
Group by $Q,P\bmod Q,D_j$. Relative to one member $D^*$, the vector
$(D-D^*)/Q$ is integral and vanishes at $j$. Its norm is $<2r$ by projection.
Its positive and negative exponents encode an ordered pair of positive integers of product
$\le\exp(2r)$, so the harmonic sum gives at most $\exp(2r)(1+2r)$ choices.
For fixed $D,Q$, the congruence $(q_1/q_2)^Q\beta(D)\equiv\pm1\bmod p$
allows at most $2Q$ residue classes of $q_1/q_2$. Each determines the reduced
fraction since $M^2<p$. Thus at most $O_{k_0,W}((1+x)(1+2r)\exp(2r))
=\exp(o(x))$ pairs fail for $W>0$ and none for $W=0$.
Choose a surviving pair. The $q_1$-coordinate of a lattice vector in the
plane forces its coefficient of $u$ to be integral. Subtraction and the line
hypothesis give the assertion for $v$. The signed character restricts to
the original one and is injective by unique factorization.
\end{proof}

\section{The determinant}\label{sec:determinant}
\begin{proposition}\label{lem:minor}
Let $p$ be prime, $K\ge1$, $m\ge0$ integers and $N=2K$. Let
$\gamma:\Z^2\to\Q^\times$ be a homomorphism into the principal $p$-units
with $\gamma(1,0)\equiv1\bmod p^m$ and, if $m>0$, $\gamma(1,0)\ne1$.
Let $F\subset\Z^2$ be finite with at least $N-1$ points and pairwise
distinct second coordinates. For $w$ with $\gamma(w)\ne1$, put
$F'=F\cup(F+w)$. At $(i,j)\in F'$ take the two blocks
$(\binom jk)_{0\le k<K}$ and $(\gamma(i,j)\binom jk)_{0\le k<K}$,
clearing the row by the positive denominator of $\gamma(i,j)$.
Suppose $|\binom jk|\le B$ throughout and every $N$ distinct sites in
$F'$ have sum of heights at most $U$. There is a nonzero integer
$N$-row minor $\Delta$ satisfying
\begin{equation}\label{eq:both}
 \begin{split}
 \log|\Delta|&\le U+\log(N!)+N\log B,\\
 4K\le3m\quad&\Longrightarrow\quad \log|\Delta|\ge\tfrac29 N^2\log p .
 \end{split}
\end{equation}
\end{proposition}
\begin{proof}
Rank failure gives a relation $P(j)+\gamma(i,j)Q(j)=0$ on $F'$ with
$P,Q$ of degree $<K$, not both zero. Since $F$ has at least $K$ distinct
second coordinates, both polynomials are nonzero. Elimination at each
point of $F$ and its translate gives $2K-1$ distinct roots of
$P(L+w_2)Q(L)-\gamma(w)P(L)Q(L+w_2)$. Its degree is at most $2K-2$ and
its leading coefficient is $(1-\gamma(w))\operatorname{lc}(P)\operatorname{lc}(Q)\ne0$.
Thus a nonzero minor $\Delta$ exists. Its row entries are bounded by
$B\exp h(\gamma(i,j))$, so determinant expansion gives the upper bound.

Work in $\Z_p$ and put $z_0=\gamma(1,0)$, $z_1=\gamma(0,1)$.
At positive depth, $z_0\ne1$ gives a finite valuation of $1-z_0$. Row denominators are
$p$-adic units. Translating the sites into the nonnegative quadrant
preserves the valuation, since binomial translation is unipotent and
the character gains a fixed $p$-unit. Modulo $p^m$ the second block has
entries $z_1^j\binom jk$. Expand $z_1^j=\sum_b\binom jb(z_1-1)^b$,
a finite sum on these rows. Counting pairs of subsets by their union
expresses $\binom jk\binom jb$ as an integral combination of
$\binom jn$ with $n\le k+b$. Thus the coefficient of $\binom jn$ in
column $k$ is divisible by $p^{\max(0,n-k)}$.

In a multilinear term with $d$ expanded columns the other $K-d$
columns contribute errors divisible by $p^m$. The first block forces
the surviving $n$'s to be distinct and at least $K$. Their sum exceeds
that of the $d$ distinct $k<K$ by at least
$\sum_{i=0}^{d-1}[(K+i)-(K-d+i)]=d^2$. Hence
$v_p(\Delta)\ge\min_{0\le d\le K}((K-d)m+d^2)$.
If $m\ge4K/3$, each term is at least
$(d-2K/3)^2+8K^2/9\ge2N^2/9$. Now $\log|\Delta|\ge v_p(\Delta)\log p$.
\end{proof}

\begin{lemma}\label{lem:diamond}
For each admissible $\eta,\sigma,\kappa$ there exists $X\ge0$ such that
every triple in $\mathcal B_X$ has a prime $p\mid s$ above $Y$ and data
as in Proposition~\ref{lem:minor} with $m=v_p(s)$ and
\begin{equation}\label{eq:budget}
           4K\le3m,\qquad U+\log(N!)+N\log B<\tfrac29 N^2\ell .
\end{equation}
\end{lemma}
\begin{proof}
Take $p,v,j$ from Proposition~\ref{lem:separate}. Since
$1-3a_0=(1-9\eta)/2>2\kappa$, the choice of $\theta$ in
Table~\ref{tab:parameters} is possible. Choose $0<Z<1$ with
\begin{equation}\label{eq:slack}
 \frac{(1+Z)(1+2Z)^2-1/3}{3(1+\theta)(1-Z)^4}<\frac29 .
\end{equation}
At $Z=0$ the left side is $2/[9(1+\theta)]$.

Apply Lemma~\ref{lem:fresh} with $h_0=h$, since $|v_j|\log 2\le h$.
For $d_0=|T|-1>0$, the arithmetic--geometric mean inequality and
$\sum_{q\in T}\log q\le3x$ give
$V_0=r^{d_0}/(d_0!\prod_{q\ne j}\log q)\ge(r/(3x))^{d_0}$.
The product of these $d_0$ primes lies between $(d_0+1)!$ and $R$, so
$(d_0+1)\log(d_0+1)-d_0\le3x$ by integrating $\log y$.
Thus $d_0\le7x/\log x$ for large $x$. With $r/(3x)=2/(3\log x)$
this gives the required volume bound, also valid when $d_0=0$ and $V_0=1$.

On the resulting plane put $\gamma(i,j')=\chi(iv+j'u)$. Then
\begin{equation}\label{eq:plane-height}
 t=\|\pi u_0\|_w+\log q_2,\qquad
 h(\gamma(i,j'))\le g(i,j'):=h|i+\xi j'|+t|j'|.
\end{equation}
Indeed the height norm is the larger of the positive and negative
weighted exponent sums. For $u-\xi v$ each sum is at most $t$, since
the new coordinates have opposite signs. The principal-unit character
$\gamma$ is injective and $\gamma(1,0)=z\ne1$ has $v_p(1-z)=m$.
We have $t>b_0\ell$ since $q_2>\sqrt M\ge\exp(\nu x/2)$.
For large $x$, $r=2x/\log x<\nu x$, so $t<\nu x+\log M$. For $p\ge H$ this is at most
$\ell-(\lambda-2\nu)x\le(3\eta+2\nu)\ell$, and for $p<H$ it is
$2\nu x\le4\nu\ell$. Both bounds are $<a_0\ell$.

The $g$-ball of radius $y$ has area $Ay^2$. A nonzero $(-P,Q)$ has
length at least $h$ for $Q=0$, at least $k_0\ell$ for $1\le|Q|\le W$
by the slope inequality, and greater than $Wb_0\ell$ otherwise.
Each exceeds $2ht/\delta=2\ell/(Z\alpha)$ for large $x$, since
$k_0,Wb_0>2/(Z\alpha)$ and $h\ge\log b\ge H-\log2$ gives $h/\ell\to\infty$ uniformly.
Minkowski's second theorem bounds the product of the successive minima
by $4/A=2ht$, so the second minimum is at most $\delta$.
Proposition~\ref{lem:separate}'s basis argument gives a basis in $g\le\delta$.

A centered unit-area cell of this basis lies in $g\le\delta$.
Translating cells gives the bounds for $C(y)=\#\{n\in\Z^2:g(n)\le y\}$,
\begin{equation}\label{eq:count}
           A(y-\delta)_+^2\le C(y)\le A(y+\delta)^2 .
\end{equation}
Take $F=\{g\le\rho\}\cap\Z^2$ and a basis vector $w$. Then
$\gamma(w)\ne1$ and $|F|\ge(1-Z)^2N_0\ge2K-1$.
Two distinct sites of $F$ with the same second coordinate would differ
by a vector of length at least $h>2\rho$, since $2\alpha a_0<1$.
Further, $N_0=2\alpha^2ht/\ell^2\to\infty$ uniformly. Since $h\le H$,
\[
 2N_0\ell\le4\alpha^2a_0 h<3(1-2\kappa)h\le3m\ell .
\]
The strict coefficient gap and $N=(1-Z)^2N_0+O(1)$ give $4K\le3m$.

All sites of $F'$ lie in $g\le\rho_1$, so $|j'|\le\rho_1/t$ and
$|\binom{j'}k|\le B$, using $|\binom{-n}k|=\binom{n+k-1}k$ for $n\ge1$.
By \eqref{eq:plane-height} each $N$-row height sum is at most
\[
 \begin{aligned}
 \sum_{g(n)\le\rho_1}g(n)
 &=\int_0^{\rho_1}\bigl(C(\rho_1)-C(y)\bigr)\,dy\\
 &\le A\bigl[\rho_1(\rho_1+\delta)^2-(\rho_1-\delta)^3/3\bigr]=:U
 \end{aligned}
\]
by \eqref{eq:count}. Since $\rho_1=(1+Z)\rho$ and $N_0\ell=3(1+\theta)\rho$,
$U/(N^2\ell)$ equals the left side of \eqref{eq:slack} plus $o(1)$ uniformly.
Also $t\ge b_0\ell$ gives $\rho_1/t=O(N)$, so $\log B=O(N)$.
Since $\log N\le N$, also $\log(N!)\le N^2$. The remaining normalized cost is
$O(1/\ell)$, which fits the strict margin in \eqref{eq:slack}.
\end{proof}

\begin{lemma}\label{lem:empty}
Assume the conclusion of Lemma~\ref{lem:diamond}. For each admissible
$\eta,\sigma,\kappa$ there exists $X\ge0$ with $\mathcal B_X=\varnothing$.
\end{lemma}
\begin{proof}
At the threshold of Lemma~\ref{lem:diamond}, a triple in $\mathcal B_X$
would yield a minor with $\tfrac29N^2\ell\le\log|\Delta|<\tfrac29N^2\ell$
by \eqref{eq:both} and \eqref{eq:budget}.
\end{proof}

\section{Proof of Theorem~\ref{thm:main}}\label{sec:end}
\begin{proposition}\label{lem:tail}
Suppose each $\eta$ admits a pair $\sigma,\kappa$ in
Table~\ref{tab:parameters} with the following property. Both exterior
bounds hold for ordered triples satisfying \eqref{eq:high} for sufficiently
large $x$. Suppose also that for each admissible $\eta,\sigma,\kappa$ some
$\mathcal B_X$ is empty. Then for each such $\eta$ there is a threshold
$X_\eta$ with
\[
                 H<R^{1/3-\eta}\qquad (x\ge X_\eta)
\]
for every triple.
\end{proposition}
\begin{proof}
Fix the pair in the first hypothesis and take the larger of the exterior
and exclusion thresholds for that pair. Exchanging $a,b$ preserves $x,H$.
An ordered triple above this threshold satisfying \eqref{eq:high} would
belong to the empty set $\mathcal B_X$.
\end{proof}
\begin{lemma}\label{lem:assembly}
A uniform polynomial bound and the two hypotheses of
Proposition~\ref{lem:tail} imply Theorem~\ref{thm:main}.
\end{lemma}
\begin{proof}
Given $\eps>0$, choose $\eta$ as in Section~\ref{sec:cap} with
$1/9-\eta<\eps$. For $x\ge X_\eta$ the last proposition gives
$H<R^{1/3-\eta}\le R^{2/9+\eps}$. For $x<X_\eta$ use
$H\le C_0R^k\le C_0\exp(3kX_\eta)$. Thus one may take
$C_\eps=\max(1,C_0\exp(3kX_\eta))$.
The exterior thresholds use the constants of Fact~\ref{fact:local}.
The polynomial bound supplies $(C_0,k)$ independently.
\end{proof}
\begin{proof}[Proof of Theorem~\ref{thm:main}]
Proposition~\ref{lem:cap} gives the polynomial bound. Use the pair
in Proposition~\ref{lem:small} and apply Proposition~\ref{lem:real}
at this same $\kappa$. Lemma~\ref{lem:empty} gives the exclusion.
Lemma~\ref{lem:assembly} concludes the proof.
\end{proof}

\section*{Use of AI tools}

The author devised the method of proof. Codex (OpenAI, GPT-6 Astra) and
Antigravity (Google, Gemini 4 Neon) ran blind reviews of the manuscript,
and Codex assisted with the Lean formalization. Claude Code (Anthropic,
Claude Fable 5) checked the proofs and the references. The author checked
every step of the proof and every reference, and takes full responsibility
for the content of the paper.

\section*{Statements and declarations}

\noindent\textbf{Competing interests.} The author has no competing
interests.

\noindent\textbf{Funding.} No funding was received for this work.

\noindent\textbf{Data availability.} A Lean 4 formalization of
Theorem~1.1 is available at \url{https://github.com/ryzfsu/abc-log-bound-two-ninths}.

\bibliographystyle{amsplain}
\bibliography{refs}
\end{document}
